% ==============================================================================
% Chapter 1: Introduction
% ==============================================================================
\pagenumbering{arabic}

\chapter{Introduction}

\section{Motivations and Importance}
I have wanted to become a game developer since I began my software engineering major. However, I never had the chance to start, because I always took web development jobs during my summer breaks. I began learning game development about one year ago. Although that is a full year on the calendar, the time I actually spent developing games was about 4 months and 2 weeks, including my previous cooperative education project in Japan. This internship adds another 3 months, so I have roughly 7 months of real game development experience in total. It is my first opportunity to learn the professional game development workflow and to follow my dream.

\section{Objectives}
\begin{enumerate}[label=\arabic*.]
    \item To learn and take part in the professional game development workflow, from pitching and documentation (GDD and TDD) through implementation and pull request review.
    \item To improve my ability to plan game architecture, including understanding the trade-offs between design choices and judging whether a design suits the game.
    \item To write cleaner and more maintainable code that meets industry standards, using code review feedback.
    \item To learn to define project scope on my own and manage time realistically, so that I can later balance a job with solo indie development.
    \item To develop professional skills such as problem-solving, presenting, and communicating ideas clearly within a team.
\end{enumerate}

\newpage
\section{Scope of Work}
This project consists of two main components: participation in the studio's professional game development workflow, and a comparison of that workflow with my previous solo development experience.

The work was carried out on a personal 2D top-down action game in Unity, which the company allowed me to develop during the internship. Following the studio's process, the work covered pre-production (pitching the game idea, writing the Game Design Document and Technical Design Document, and designing the UML class diagram), followed by a prototype phase in which core mechanics were built as a proof of concept. The prototype was intended to validate the design before moving into production toward a Minimum Viable Product (MVP). During production, all code was submitted through pull requests and reviewed by the development team under the studio's standards for naming, structure, and design.

Out of scope for this project: completing the MVP or the full game, final art, audio, and polish, a release build, and marketing. The internship was scoped as an opportunity to learn the professional workflow, not to deliver a finished product.
\section{History and Detail of Company}
\begin{itemize}[label=-]
    \item \textbf{Name and location}: \companyName, Bangkok, Thailand.
    \item \textbf{Company profile; products or main services of the company}: A game development company that develops video games across multiple genres using Unity.
    \item \textbf{Organization and management model}: The company follows an Agile development process based on Scrum. Work is organized into sprints and the team holds a daily morning meeting along with a separate meeting for engineers. All code changes go through pull request review by team members before merging.    
    \item \textbf{Position or work detail assigned to student}: Game Developer Intern, responsible for a personal 2D top-down action game in Unity.
    \item \textbf{Name and position of student's supervisor}: \supervisorName, \supervisorPosition.
    \item \textbf{Internship/cooperative education period}: \coopStartDate{} -- \coopEndDate.
\end{itemize}

\section{Expected Benefits}
This project is expected to provide benefits in the following areas:

\begin{itemize}[label=-]
    \item \textbf{To myself}: Practical experience of the professional game development workflow, including pitching, documentation (GDD, TDD, UML), code review, and sprint-based work. It also helps me improve my architecture planning, write more maintainable code, and understand the trade-offs behind design decisions, which I can apply to my future career and my own indie projects.
    \item \textbf{To the organization}: An opportunity to evaluate and mentor a prospective junior developer.
    \item \textbf{To the university and academic field}: A documented case study of a student moving from solo development to a professional team workflow, including the time spent on each phase and the code review process. This may serve as a reference for future cooperative education students and for the connection between the university and the game industry.
\end{itemize}